\subsection{Задача Гурса}
\begin{theorem}[существования и единственности]
	Решение $u(x, y) \in C^2(\Pi) \cap C^1(\overline{\Pi})$ задачи Гурса:
	\[
		Lu = u_{xy} + a u_x + b u_y + cu = f(x, y),\ (x, y) \in \Pi = (0, l) \times (0, h) \quad (1.f)
	\]
	\[
		\textbf{Условия на характеристиках(УХ): } u|_{x=0} = \phi(y),\ 0 \le y
		\le h,\quad
		u|_{y=0} = \psi(x),\ 0 \le x \le l
	\]
	\[
		\textbf{Условия согласования(УС): } \phi(0) = \psi(0) = u(0, 0),
	\]
	\[\text{если } a(x, y) \in C^1(\overline{\Pi}),\ b(x, y), c(x, y), f(x, y) \in C^1(\overline{\Pi}), \phi(y) \in C^2([0, h]),\ \psi \in C^2([0, l]),
	\] существует и единственно.
\end{theorem}

Физический смысл: "сопло Лаваля".

\subsection{Смешанная задача(СЗ) для однородного уравнения колебаний полубесконечной струны(ОУКПБС), закрепленной на левом конце.}

\begin{align*}
	                       & u_{tt} = a^2 u_{xx}, \quad x > 0,\ t > 0 \tag{1.0}    \\
	\textbf{НУ: } \quad    & u|_{t=0}=u_0(x),\ u_t|_{t=0}=u_1(x),\ x \ge 0 \tag{2} \\
	\textbf{ОГУ1Т: } \quad & u|_{x=0}=0,\ t \ge 0 \tag{3.1.0}                      \\
	\textbf{УС: } \quad    & u_0(0)=u_1(0)=0=u_0^{''}(0), \tag{4.1.0}
\end{align*}

$u(x, t) \in C^2((0, +\infty) \times (0, \infty)) \cap C^1([0, +\infty) \times [0, \infty))$

$u_0(x) \in C^2(0, +\infty),\quad u_1(x) \in C^1(0, +\infty)$

Вспомним Задачу Коши(ЗК) для ОУКБС:
\begin{align*}
	                    & u_{tt} = a^2 u_{xx}, \quad x > 0,\ t > 0 \tag{1.0}             \\
	\textbf{НУ: } \quad & u|_{t=0}=u_0(x),\ u_t|_{t=0}=u_1(x),\ x \in \mathbb{R} \tag{2}
\end{align*}
$u(x, t) \in C^2(t>0) \cap C^1(t \ge 0),\quad u_0(x) \in C^2(\mathbb{R}),\quad u_1 \in C^1(\mathbb{R})$

\begin{lemma}
	Если НУ в ЗК (1.0), (2) являются нечётными функциями относительно т. $x=x_0$, то соответствующее решение $u(x, t)$ ЗК (1.0), (2) равно нулю в т. $x=x_0$, т.е. $u(x_0, t) = 0,\quad t \ge 0$.
\end{lemma}

\begin{lemma}
	Если НУ в ЗК (1.0), (2) являются чётными функциями относительно т. $x=x_0$, то частная производная соответствующего решения $u(x, t)$ ЗК (1.0), (2) равна нулю в т. $x=x_0$, т.е. $u_x(x_0, t) = 0,\quad t \ge 0$.
\end{lemma}

\begin{note}
	функция $y = f(x)$, определенная на $(a, b)$ и симметричная относительно точки $x=x_0 \in (a, b)$ называется чётной(нечётной), если $f(2x_0 - x) = f(x)(f(2x_0 - x) = -f(x))$, где $f(x_0 - x^{'}) = \pm f(x_0 + x^{'})$(относительно т.$x_0$).
\end{note}

\begin{proof}
	(Доказательство первой леммы)
	Положим $x_0 = 0$. Тогда по теореме о существовании и единственности решения ЗК (1.0), (2) для ОУКБС в рассматриваемых классах функций $\exists!$ и представляется Ф.Д.0(формулой Даламбера):
	\[
		u(x, t) = \frac{1}{2} [u_0(x-at) + u_0(x+at)] + \frac{1}{2a}\int_{x-at}^{x+at}u_1(\xi)d\xi \Rightarrow
	\]
	\[
		u(0, t) = \frac{u_0(-at)+u_0(at)}{2} + \frac{1}{2a} \int_{-at}^{at} u_1(\xi)d\xi = 0,
	\]
	при $t \ge 0$, т.к. из условий леммы функции $u_0(x)$ и $u_1(x)$ --- нечётные.

	Аналогично, доказывается вторая лемма с учётом того факта, что если $f(x)$ --- чётная функция, то её производная $f^{'}(x)$ --- нечётная.
\end{proof}

\subsection{Метод продолжений}

\textbf{Следствие 1} Для того, чтобы свести СЗ (1.0), (2), (3.1.0), (4.1.0) для ОУКПБС, закрепленной на левом конце, к ЗК для ОУКБС нужно начальные функции $u_0(x)$ и $u_1(x)$ нечетным образом продолжить влево за т. $x=0$.

\textbf{Следствие 2} Для того, чтобы свести СЗ для ОУКПБС, со свободным левым концом, к ЗК для ОУКБС нужно начальные функции $u_0(x)$ и $u_1(x)$ четным образом продолжить влево за т. $x=0$.

\begin{theorem}
	Решение $u(x, t)$ СЗ (1.0), (2), (3.1.0), (4.1.0) для ОУКПБС, закрепленной на левом конце, существует и единственно и представляется обобщенной Ф.Д.0(Ф.Д.0.1):
\end{theorem}

\[
	u(x,t)=
	\begin{cases}
		\dfrac{u_0(x-at)+u_0(x+at)}{2}+\dfrac{1}{2a}\displaystyle\int_{x-at}^{at+x} u_1(\xi)\,d\xi, & x\ge at,\ t\ge 0, \\[1.2em]
		\dfrac{u_0(at+x)-u_0(at-x)}{2}+\dfrac{1}{2a}\displaystyle\int_{at-x}^{at+x} u_1(\xi)\,d\xi, & 0<x<at,\ t\ge 0.
	\end{cases}
\]

Получим данное решение.
\begin{proof}
	В соответствии со \textbf{Следствием 1}, сведем эту СЗ к ЗК для ОУКБС, продолжив нечетным образом НУ влево за т. $x=0$. Таким образом, получим следующую ЗК:

	\begin{align*}
		                    & u_{tt} = a^2 u_{xx}, \quad x \in \mathbb{R},\ t > 0 \tag{1.0} \\
		\textbf{НУ: } \quad & u|_{t=0}=U_0(x)=
		\begin{cases}
			u_0(x),\ x \ge 0 \\
			-u_0(-x),\ x < 0
		\end{cases}
		,\ u_t|_{t=0}=U_1(x) =
		\begin{cases}
			u_1(x),\ x \ge 0 \\
			-u_1(-x),\ x < 0
		\end{cases},
		\ x \in \mathbb{R} \tag{2}
	\end{align*}

	Решение $u(x, t)$ ЗК для ОУКБС существует и представляется Ф.Д.0:
	\[
		u(x, t) = \frac{U_0(x-at) + U_0(x+at)}{2} + \frac{1}{2a} \int_{x-at}^{x+at} U_1(\xi) d\xi,\ x \in \mathbb{R},\ t > 0
	\]
	При $x \ge at,\ t \ge 0\ (x-at \ge 0)$:
	\[
		u(x, t) = \frac{u_0(x-at) + u_0(x+at)}{2} + \frac{1}{2a} \int_{x-at}^{at+x} u_1(\xi) d\xi,\ x \ge at,\ t \ge 0
	\]
	При $0 \le x \le at,\ t \ge 0\ (x-at \le 0)$:
	\[
		u(x, t) = \frac{u_0(at+x) - u_0(at-x)}{2} + \frac{1}{2a} \int_{at-x}^{at+x} u_1(\xi) d\xi,\ 0 < x < at,\ t \ge 0
	\]
\end{proof}

Физический смысл: ОР ОУКПБС: $u(x, t) = f(x-at) + g(x+at)$.

Из условия закрепления левого конца:
\[
	u|_{x=0} = f(-at) + g(at) = 0,\ t \ge 0 \Rightarrow \\
	f(s) = -g(s),\ s=-at \le 0
\]
Т.о., при $0 < x < at \Rightarrow u(x, t) = g(at + x) - g(at - x)$.

Т.о., наблюдается явление отражения волны, бегущей справа налево от левого(закрепленного) конца струны с сохранением абсолютной величины и изменением знака на противоположный.

СЗ для ОУКПБС, со свободным левым концом
\begin{align*}
	                       & u_{tt} = a^2 u_{xx}, \quad x > 0,\ t > 0 \tag{1.0}    \\
	\textbf{НУ: } \quad    & u|_{t=0}=u_0(x),\ u_t|_{t=0}=u_1(x),\ x \ge 0 \tag{2} \\
	\textbf{ОГУ2Т: } \quad & u_x|_{x=0}=0,\ t \ge 0 \tag{3.2.0}                    \\
	\textbf{УС: } \quad    & u_0^{'}(0)=u_1(0)=0, \tag{4.2.0}
\end{align*}

$u(x, t) \in C^2((0, +\infty) \times (0, \infty)) \cap C^1([0, +\infty) \times [0, \infty))$

$u_0(x) \in C^2[0, +\infty),\quad u_1(x) \in C^1[0, +\infty)$

\begin{theorem}[о существовании и единственности]
	Решение $u(x, t)$ СЗ $(1.0),\ (2),\ (3.2.0),\ (4.2.0)$ для ОУКПБС со свободным левым концом в рассматриваемом классе функций существует и единственно и представляется Ф.Д.0.2:
	\[
		u(x,t)=
		\begin{cases}
			\dfrac{u_0(x-at)+u_0(x+at)}{2}+\dfrac{1}{2a}\displaystyle\int_{x-at}^{x+at} u_1(\xi)\,d\xi,                                                       & x\ge at,\ t>0,   \\[1.2em]
			\dfrac{u_0(x+at)+u_0(x-at)}{2}+\dfrac{1}{2a}\left(\displaystyle\int_{0}^{at-x} u_1(\xi)\,d\xi+\displaystyle\int_{0}^{at+x} u_1(\xi)\,d\xi\right), & 0\le x<at,\ t>0.
		\end{cases}
	\]
\end{theorem}

Физический смысл: ОР ОУКПБС: $u(x, t) = f(x-at) + g(x+at);$

Из условий свободности левого конца:
\[
	u_x|_{x=0} = f^{'}(-at) + g^{'}(at),\ t \ge 0 \Rightarrow f(s) = -g(-s) + C,\ s = -at \le 0
\]
Но т.к. $u_t|_{x=0} = f^{'}(s) - g^{'}(-s) = 0$, то $f(s) = g(-s) - C$

$u(x, t) = g(at + x) + g(at - x) - C$ в области $0 < x < at$.

Т.о., наблюдается явление отражения волны, бегущей справа налево от левого(свободного) конца струны с сохранением абсолютной величины и знака.